\chapter[Pullback Metric]{Pullback Riemannian Metric}\label{ch:sec:pullback-riemannian-metric}
Throughout this chapter let $F:\mathcal P \longrightarrow F(\mathcal P)$ denote the global diffeomorphism established in Chapter~\ref{ch:sec:local-diffeomorphism}.
The image \(F(\mathcal P)\) is an open subset of the finite-dimensional vector space
\[
\prod_{c\in\mathcal C}\operatorname{Im}(Q_c),
\]
which is naturally endowed with the Euclidean inner product inherited from the ambient spaces.
The purpose of this chapter is to transport this Euclidean geometry back to the posterior manifold.

The pullback metric is the unique Riemannian structure that makes the logarithmic map a local isometry. This metric defines geodesic distances on the posterior manifold, which in turn determine the Lipschitz constants for information functionals (Chapters~13--16) and the variational structure of the UOD (Chapters~20--25).

\section{Euclidean Metric on the Image}\label{ch:sec:euclidean-metric-on-the-image}
For every point
$y\in F(\mathcal P),$
define the Euclidean metric
$g_E{}_y(u,v)=\langle u,v\rangle,$
where \(\langle\cdot,\cdot\rangle\) denotes the standard Euclidean inner product.
Since the ambient space is linear, this metric is smooth.
\begin{proposition}[Smoothness of the Euclidean Metric]
\label{ch:prop:7-1}
The Euclidean metric defines a smooth Riemannian metric on \(F(\mathcal P)\).
\end{proposition}
\begin{proof}
The Euclidean inner product is constant in every affine coordinate system. Hence it varies smoothly and is positive definite on every tangent space. 
\end{proof}
\paragraph{Why the pullback metric?}
The image $F(\mathcal P)$ is a subset of Euclidean space, so it inherits the Euclidean metric.  But what we need is a metric on the posterior manifold itself---a way to measure how much two distributions differ in terms of their cryptographic structure.  The pullback metric $g=F^*g_E$ is the unique Riemannian metric that makes $F$ a local isometry: the distance between $P$ and $Q$ on the manifold equals the Euclidean distance between their images $F(P)$ and $F(Q)$.  This identification is the bridge between geometry (distributions) and algebra (WIS factorisation).

\section{Pullback Metric}\label{ch:sec:pullback-metric}
\begin{definition}[Pullback Metric]
\label{ch:def:7-2}
The pullback metric induced by the canonical logarithmic coordinate map is
\[
g
=
F^{*}g_E.
\]
Equivalently, for every
$P\in\mathcal P$
and tangent vectors
$u,v\in T_P\mathcal P,$
define
\[
g_P(u,v)
=
g_E{}_{F(P)}
\!\left(
DF_P(u),
DF_P(v)
\right).
\]
\end{definition}

\paragraph{Why this metric.}
The pullback metric is not chosen by the analyst; it is forced by the logarithmic representation.  Unlike the Fisher metric, which is selected for its statistical properties, this metric is intrinsic to the encryption system.

\begin{proposition}[Well-Definedness]
\label{ch:prop:7-3}
The pullback metric is well defined.
\end{proposition}
\begin{proof}
By Theorem~\ref{ch:thm:5-4}, \(F\) is a diffeomorphism onto its image. Therefore the differential
$DF_P:T_P\mathcal P\rightarrow T_{F(P)}F(\mathcal P)$
is a linear isomorphism for every \(P\).
Since \(g_E\) is defined on every tangent space of the image, the expression above is well defined. 
\end{proof}
\begin{proposition}[Symmetry]
\label{ch:prop:7-4}
For every posterior point \(P\), the bilinear form
$g_P$
is symmetric.
\end{proposition}
\begin{proof}
Because
$g_E(u,v)=g_E(v,u),$
we obtain
\[
g_P(u,v)
=
g_E(DF_Pu,DF_Pv)
=
g_E(DF_Pv,DF_Pu)
=
g_P(v,u). \qquad
\]
\end{proof}
\begin{proposition}[Positive Definiteness]
\label{ch:prop:7-5}
The pullback metric is positive definite.
\end{proposition}
\begin{proof}
Let
$u\neq0.$
Since
$DF_P$
is an isomorphism,
$DF_P(u)\neq0.$
Because the Euclidean metric is positive definite,
$g_E(DF_Pu,DF_Pu)>0.$
Hence
$g_P(u,u)>0.$
Therefore \(g\) is positive definite. 
\end{proof}
\begin{theorem}[Smooth Riemannian Metric]
\label{ch:thm:7-6}
The pullback metric defines a smooth Riemannian metric on the posterior manifold.
\end{theorem}
\begin{proof}
By Proposition~\ref{ch:prop:7-1} the Euclidean metric is smooth.
The pullback of a smooth tensor field by a smooth map is smooth.
Propositions~\ref{ch:prop:7-4} and~\ref{ch:prop:7-5} establish symmetry and positive definiteness.
Hence \(g\) is a smooth Riemannian metric on \(\mathcal P\). 
\end{proof}
\section{Coordinate Representation}\label{ch:sec:coordinate-representation}
Let
$J_P=DF_P.$
The pullback metric satisfies
\[
g_P
=
J_P^{\!T}J_P,
\]
where the transpose is taken with respect to the Euclidean inner product.
This identity follows immediately from the definition of the pullback metric.
\begin{corollary}[Positive Definite Matrix Representation]
\label{ch:cor:7-7}
The matrix representing the pullback metric is positive definite at every posterior point.
\end{corollary}
\begin{proof}
Since
$J_P$
is invertible,
$J_P^{\!T}J_P$
is symmetric positive definite by elementary linear algebra. 
\end{proof}
\subsection{Geometric Interpretation}\label{ch:sec:geometric-interpretation}
The pullback metric measures infinitesimal posterior variations after eliminating the additive gauge symmetry through the quotient projection.
Consequently, distances, norms, gradients and quadratic forms computed with respect to \(g\) are intrinsic properties of the posterior manifold and do not depend on the choice of logarithmic representative.
\paragraph{Running example: differential privacy.}
For differential privacy (Section~\ref{ch:sec:differential-privacy-and-database-queries}), the pullback metric $g=F^*g_E$ measures the cost of moving between neighbouring output distributions.  For the Laplace mechanism with $\mathcal D=6/\varepsilon^2$, the pullback metric at $P$ acts on tangent vectors as $g_P(u,v)=\langle Q\,dL_P(u),\,Q\,dL_P(v)\rangle$, where $dL_P(u)=(u_i/P_i)_i$ is the differential of the logarithmic embedding and $Q$ projects onto the quotient.  Near the uniform prior $U_n$ the entries of $dL_P$ are of order $n$, and a direct computation gives $g_P(u,v)\approx (\varepsilon^2/2)\langle u,v\rangle$---a direct link between the privacy budget and the geometry of the posterior manifold.

This metric provides the geometric foundation for the Universal Optimality Defect introduced in the next chapter.
